\documentclass[12pt,a4paper]{article}
\usepackage[utf8]{utf8}
\usepackage{geometry}
\geometry{margin=1in}
\usepackage{hyperref}
\usepackage{listings}
\usepackage{booktabs}

\title{\textbf{ULK POLYTECHNIC INSTITUTE}\\[0.5em]\Large Cryptography \& Network Security Technical Report}
\author{Likambo Kulang John}
\date{\today}

\begin{document}

\maketitle

\section{Risk Assessment}
\subsection{Assets, Vulnerabilities, and Consequences}
A security assessment of the polytechnic institute's infrastructure revealed critical security vulnerabilities across three main assets:
\begin{itemize}
    \item \textbf{Student Records Server:} Exposed to unencrypted file transfers (FTP/HTTP) and guest network access, creating a high risk of sensitive student record interception and privacy violations.
    \item \textbf{Staff Credentials:} Vulnerable to weak password policies, making accounts susceptible to brute-force attacks and unauthorized administrative access.
    \item \textbf{Server Operating System:} Running outdated software packages, exposing the server to known exploits and potential remote code execution (RCE).
\end{itemize}

\subsection{Risk Rankings and Recommended Controls}
\begin{table}[h!]
\centering
\begin{tabular}{@{}llll@{}}
\toprule
\textbf{Risk} & \textbf{Likelihood} & \textbf{Impact} & \textbf{Recommended Control} \\ \midrule
Guest Network Access & High & High & Network Segmentation \& Firewall Filtering \\
Weak Staff Passwords & High & Medium & Strong Password Policy \& MFA \\
Unencrypted Transfers & Medium & High & Enforce SFTP / TLS Encryption \\ \bottomrule
\end{tabular}
\end{table}

\section{Encryption and Integrity Application}
A Python security toolkit (\texttt{toolkit.py}) was developed using the \texttt{cryptography} library to manage file confidentiality and integrity:
\begin{itemize}
    \item \textbf{Encryption \& Decryption:} Symmetric Fernet encryption was applied to \texttt{sample\_record.csv}. Decryption tests confirmed exact restoration of original record content.
    \item \textbf{Integrity Verification:} SHA-256 cryptographic hashing was implemented. Verification tests proved that altered files trigger immediate integrity alerts, while untampered files verify successfully.
    \item \textbf{Key Safety:} The encryption key (\texttt{secret.key}) is stored locally and excluded from version control via \texttt{.gitignore}.
\end{itemize}

\section{Network Traffic Filtering}
Traffic filtering was configured using \texttt{iptables} rules to isolate the server:
\begin{verbatim}
sudo iptables -F
sudo iptables -A INPUT -m conntrack --ctstate ESTABLISHED,RELATED -j ACCEPT
sudo iptables -A INPUT -s 192.168.20.0/24 -j DROP
sudo iptables -A INPUT -s 192.168.10.0/24 -p tcp --dport 22 -j ACCEPT
sudo iptables -A INPUT -p tcp --dport 22 -j DROP
\end{verbatim}
Testing confirmed that connection requests from the Staff network succeeded, while guest network access and external connection attempts timed out as expected.

\section{Conclusion and Submission}
The technical controls effectively mitigate unauthorized access, safeguard student data in transit, and isolate sensitive server resources.

\textbf{GitHub Repository URL:} \url{https://github.com/Moxi69-beep/cryptography-network-security}

\end{document}